A robótica móvel aplicada a ambientes internos exige que percepção, controle, localização e mapeamento atuem de forma integrada. Em plataformas comerciais de maior custo, parte dessas demandas pode ser atendida por sensores dedicados, unidades de processamento mais robustas e sistemas de navegação consolidados. Em protótipos de baixo custo, porém, a viabilidade da solução depende de escolhas técnicas mais restritas: sensores acessíveis, processamento embarcado limitado, comunicação confiável entre componentes e uma arquitetura de software capaz de organizar dados heterogêneos em tempo real.

Nesse contexto, o SLAM (\textit{Simultaneous Localization and Mapping}) ocupa papel central, pois permite que o robô estime sua própria pose enquanto constrói uma representação do ambiente. Sensores RGB-D, como o Kinect, oferecem uma alternativa economicamente acessível a sensores LiDAR ao combinar imagem e profundidade em um único dispositivo \cite{jin2019visualslamrgbd, lemos2019plataformaroboticamovel}. Entretanto, o uso de informação visual e de profundidade aumenta a demanda por processamento, o que torna a execução embarcada em placas como o Raspberry Pi uma questão técnica relevante \cite{raspberrypi2019raspberrypi, balado2025systematicliteraturereview}.

Este trabalho investiga essa questão a partir de um robô físico de baixo custo construído com o reaproveitamento do chassi de um aspirador comercial. A plataforma integra um Raspberry Pi 4, uma câmera RGB-D Kinect v1, um Arduino Nano, atuadores, sensores auxiliares e uma pilha de software baseada em ROS 2 Jazzy \cite{arduino2024arduinonano, quigley2009rosopensource, openrobotics2024jazzy}. A proposta não se restringe à montagem mecânica do protótipo: o foco está na organização e avaliação do pipeline necessário para adquirir dados sensoriais sincronizados, publicar odometria e transformações, executar o RTAB-Map e registrar evidências que permitam avaliar a viabilidade do mapeamento em tempo real.

O escopo deste Trabalho de Conclusão de Curso II prioriza a validação no robô físico em ambiente real. Simulação em Gazebo e processamento remoto, embora relevantes como alternativas de desenvolvimento, deixam de ser eixos de avaliação para dar lugar à análise das restrições impostas pelo hardware embarcado. Da mesma forma, a navegação completa com Nav2 é tratada como continuidade futura, condicionada à estabilização prévia da localização e do mapa. Assim, a contribuição deste trabalho é consolidar uma estrutura técnica e metodológica que permita discutir a viabilidade do SLAM visual embarcado em uma plataforma acessível, enfrentando desafios como a instabilidade na aquisição simultânea de dados RGB e profundidade.

A pesquisa adota a metodologia \textit{Design Science Research} (DSR), adequada para trabalhos que constroem e avaliam artefatos tecnológicos em contexto aplicado \cite{hevner2004designscienceinformation, dresch2015distinctiveanalysiscase}. O artefato deste trabalho é a plataforma robótica integrada, composta pelo hardware físico, pelos sensores, pela interface de controle de baixo nível e pelo conjunto de nós ROS 2 responsáveis pela aquisição, sincronização e processamento dos dados. A avaliação enfatiza a delimitação do pipeline, a organização do protocolo experimental e a análise das evidências coletadas durante os testes de mapeamento.

Diante disso, a questão de pesquisa é formulada da seguinte forma:

\begin{center}
    \textit{Como estruturar e avaliar um pipeline de SLAM visual embarcado em um robô móvel físico de baixo custo, considerando as limitações de sensores acessíveis, odometria, transformações e processamento em hardware restrito?}
\end{center}

\section{Objetivos}

O objetivo geral deste trabalho é estruturar, integrar e avaliar progressivamente uma plataforma robótica móvel de baixo custo para execução de SLAM visual embarcado em ambiente interno, tendo como foco a estabilidade do pipeline de sensores, odometria, transformações e mapeamento no robô físico.

Para alcançar esse objetivo geral, foram definidos os seguintes objetivos específicos:

\begin{enumerate}
    \item fundamentar os conceitos de robótica móvel, SLAM visual, sensores RGB-D, ROS 2, odometria e fusão sensorial aplicados ao protótipo;
    \item descrever a arquitetura física e lógica do robô, incluindo Raspberry Pi, câmera RGB-D, Arduino Nano, atuadores, sensores auxiliares e comunicação entre componentes;
    \item configurar a pilha de software necessária para publicação de tópicos, odometria, transformações de coordenadas e execução do RTAB-Map no robô físico;
    \item definir um protocolo experimental para avaliação do SLAM embarcado em testes reais controlados;
    \item preparar a estrutura de análise dos resultados, contemplando estabilidade da pose, coerência da odometria, qualidade do mapa, taxa de processamento e efeitos de ajustes de fusão sensorial.
\end{enumerate}

\section{Organização do Trabalho}

Este documento está estruturado da seguinte forma. O Capítulo \ref{cap-fundamentacao-teorica} apresenta a fundamentação teórica relacionada a robótica móvel, SLAM visual, sensores RGB-D, ROS 2 e execução em hardware limitado. O Capítulo \ref{cap-metodologia} descreve a classificação da pesquisa, as etapas de desenvolvimento do artefato e o protocolo de avaliação previsto. O Capítulo \ref{cap-trabalhos-relacionados} discute trabalhos correlatos sobre SLAM, plataformas robóticas e avaliação em sistemas embarcados. O Capítulo \ref{cap-desenvolvimento} apresenta a arquitetura e a integração do protótipo físico. O Capítulo \ref{cap-resultados-discussao} organiza os resultados esperados, os critérios de validação e as evidências a consolidar após os testes em andamento. Por fim, o Capítulo \ref{cap-conclusao} apresenta o fechamento parcial da entrega, as limitações atuais e os próximos passos para a versão final do TCC II.
